% !PRINTED 119
\weberpage{119}
%<<CONT
bezeichnen können. Auch hier gilt das Gesetz
\weq{6}{\pi_a \pi_{a'} = \pi_{aa'},}
wenn $a'$ gleichfalls irgend ein Element aus $P$ ist. Es ist noch
die Frage, ob verschiedene Elemente $a$ dieselbe Permutation $\pi$
hervorrufen können. Wenn $\pi_a = \pi_{a'}$ ist, so folgt aus (6):
\[
\pi_{a'a^{-1}} = 1,
\]
und wenn wir also mit $a_0$ alle der Bedingung
\weq{7}{\pi_{a_0} = 1}
genügenden Elemente bezeichnen, so ist $a' = a_0 a$, und es kommt
darauf an, die Gesammtheit der Elemente $a_0$ zu ermitteln. Wenn
aber
\weq{8}{Rba_0 = Rb}
sein soll, so muss
\[
ba_0 = cb
\]
oder
\[
a_0 = b^{-1}cb
\]
sein, d.\,h. $a_0$ muss der Gruppe $b^{-1}Rb$ angehören, und dies ist
auch hinreichend für das Bestehen von (8).

Daraus ergiebt sich, dass die der Bedingung (7) genügenden
Elemente $a_0$ die ganze Gruppe $R_0$ erfüllen, die der Durchschnitt
aller mit $R$ conjugirten Theiler von $P$:
\[
R,\ b_1^{-1}Rb_1,\ b_2^{-1}Rb_2,\ \ldots,\ b_{j-1}^{-1}Rb_{j-1},
\]
und, wie wir früher nachgewiesen haben, ein Normaltheiler
von $P$ ist.

Daraus ergiebt sich für den Fall, dass $R_0 = 1$ ist, der Satz:

2.\ \gesperrt{Hat eine Gruppe $P$ einen Theiler vom Index $j$,
der mit seinen conjugirten Theilern ausser
dem Einheitselemente kein Element gemein
hat, so ist $P$ (einstufig) isomorph mit einer
transitiven Permutationsgruppe von $j$ Ziffern.}

Dass die Permutationsgruppe transitiv ist, sieht man aus (4),
wo für $a$ jedes der Elemente $b_1, b_2, \ldots, b_{j-1}$ gesetzt werden kann.

Der Satz 2. findet immer statt, wenn $P$ eine \gesperrt{einfache Gruppe}
und $R$ ein echter Theiler von $P$ ist, weil dann $R_0$, als Normaltheiler von $P$, sich auf die Einheitsgruppe reduciren muss\footnotemark.

Es mögen jetzt $R$ und $Q$ irgend zwei Theiler der Gruppe $P$
bedeuten, $R$ vom Index $j$, $Q$ vom Grade $q$.
\footnotetext{\gesperrt{Hölder}, Mathematische Annalen, Bd.\ 40, S.\ 57.}
